\documentclass[11pt]{article}
\usepackage[margin=1in]{geometry}
\usepackage{amsmath,amssymb,amsthm,hyperref,booktabs}
\newtheorem{theorem}{Theorem}
\newtheorem{lemma}{Lemma}
\newcommand{\polylog}{\operatorname{polylog}}
\title{A short inverse-hash circuit for a prescribed elliptic sum}
\author{Project research note}
\date{17 November 2026}
\begin{document}
\maketitle

Let $b=\lceil\log_2p\rceil$. Use the input class, canonical assignments,
independent symmetric abscissa hashes, and supplied group order
$n_E=\#E(\mathbb F_p)$ of the project's compressed relation-search
paper. Under its finite packet Assumption A.1 and the inherited C38
and relation-reduction results, the new leading time bound is
\[
 \log_2\mathbb E T\ \le\
 (26.2933775156+o(1))\,b/(\log_2 b)^{1/4}.
\]
The previous coefficient was $27.0138774237$. The change is in the
leading coefficient. The power of $b$ and the logarithmic saving are
unchanged. The hash identity, its local circuit, and the reduction
from the ordinary count to one guarded count are algebraic results.
The conversion of circuit size to a runtime has the same substantive
conditional foundation as checkpoint 004.

\section{Check the short preimage at the existing abscissa owners}
Put $M=2^b$, $\tau=2^\ell$. The slot retains a canonical abscissa $x<p$
when $t=(\alpha x+\beta)\bmod M$ lies in $[0,\tau)$, with $\alpha$ odd.
Put $A=\alpha^{-1}\bmod M$ and $C=-A\beta\bmod M$, both canonical.
The equivalent constraint is
\[
     0\le t<2^\ell,\qquad x=(At+C)\bmod M .
\]
Each accepted $x$ supplies exactly one $t$. The paid threshold has
$\ell=b/m+O(\log b+\log m)$. In the regime $m\to\infty$,
$m=\polylog b$, we have $\ell=o(b)$. The new circuit uses this
asymmetry while checking the original hash for every seed.

Use the original main Fourier parameters: limb radix $B=2^k$ with
$k=\Theta((\log b)^2)$, transform length $L=\Theta(b/k)$, and second
dimension $V=L_2=\Theta(\sqrt L)$, all transform dimensions powers of
two. The original source processor of residue $j$ inverts its first
transform and owns the actual canonical $x$ digits at indices
congruent to $j\bmod V$. The existing modulus $M_F$ and its principal
$L$th root $\omega$ have word width $W=2k+O(\log b)$.
Write $n=\lceil b/k\rceil$, $s=\lceil\ell/k\rceil$.
Choose a power of two $N$ in $[n/\sqrt m,2n/\sqrt m]$, rounding up
at the left endpoint, and put
\[
 H=V\left\lfloor\frac{N-s+1}{V}\right\rfloor,\qquad J=\lceil n/H\rceil .
\]
Eventually $V\mid N\mid L$, $H\ge s$, $N=o(n)$, and
$H/N=1-O((s+V)/N)=1-o(1)$. Finite inputs outside the guards can use
the original hash compiler. For $\ell=0$ the preimage is zero, so a
direct comparison with $C$ also suffices.

Split the public multiplier into blocks of $H$ limbs. For limb
polynomials this means $A(Z)=\sum_{h=0}^{J-1}Z^{hH}A_h(Z)$ with
$\deg A_h<H$. Let $U(Z)$ be the $s$-limb polynomial of the private
index $t$. Every $A_h U$ has degree at most $H+s-2<N$.
Compute all $J$ products from one forward transform of $U$, using the
length-$N$ transform over the old modulus with root $\omega^{L/N}$.
Its two dimensions are $N/V$ and $V$. Exactly as in the original
$GH$ factorization, one middle vertex per first frequency, per slot,
performs the second transform, the $J$ pointwise multiplications by
public spectra, and the $J$ inverse second transforms. There are
$N$ forward boundary words and $JN$ output boundary words per slot.
No graph words carry public spectra.

For each residue $j\bmod V$ the first-stage row can invert $U$'s
boundary to verify that all private digits are within range, with
the final limb restricted to the remaining bits and all padding zero.
It also inverts the $J$ corresponding output rows. A coefficient at
local position $d$ contributes to global position $hH+d$.
Since $V\mid H$, both positions have precisely the same owner. That
owner already reads the corresponding $x$ digit from the main C38
source transform. There is no new transport of abscissas, overlapping
coefficients, or range-test inputs. Each block coefficient is at most
$s(B-1)^2<M_F$, so its nonnegative integer value is unambiguous.

Starting at $c_0=0$, add the digit $C_i$, the relevant block
coefficients, and $c_i$. Compare the low radix digit with $x_i$ and
propagate the radix quotient $c_{i+1}$. At the last position compare
only the remaining $b-k(n-1)$ bits and discard the computed overflow.
The combined convolution coefficient has at most $s$ products even at
an overlap, since the blocks partition the positions of $A$. Induction
therefore gives $0\le c_i\le s(B-1)+1$. Canonical unsigned carries
need $k+O(\log b)$ bits.

Merge $R$ consecutive main source owners where $R$ is a power of two
of order $\log b$ dividing $V$. Each run of global digits of length
$R$ has one owner. Use initialized canonical carry ports only between
runs, with degree-two relays as needed. This costs
$O(mn(k+\log b)/R)=O(mb/\log b)$ binary edges. Existing range and
distinctness tests are retained. Old relays whose endpoints were
merged may keep two distinct edges to the same merged processor.
The graph is bipartite and loop-free. The new hash middle vertices
are opposite the merged source vertices.

\begin{lemma}
For fixed old non-hash source data, each accepted assignment has a
unique extension to the replacement hash interfaces. The old and new
accepted assignments are in bijection after dropping their respective
hash-only wires.
\end{lemma}
\begin{proof}
Row inversion recovers canonical digits of $t$. The checked middle
outputs and the two invertible transforms force the ordinary product
coefficients. Carries are initialized and determined by the recurrence,
with no final free word. Hence acceptance forces $x\equiv At+C\bmod M$
and the range forces $t=(\alpha x+\beta)\bmod M$. Conversely this
$t$, when in range, fixes all forward words, output words, intermediate
carry encodings, and relay values. The old hash outputs and their
carries were similarly deterministic given $x$. Restricting original
variables or valid new symbol pins preserves exact counting.
\end{proof}

\paragraph{Degree and construction.}
Each new middle has $O(V(J+1)W)$ Boolean ports. At each merged source
owner the added word degrees are at most $mR(N/V)(J+1)W$ and a given
new middle contributes at most $R(J+1)W$ parallel bits.
Here $J=O(\sqrt m)$ and $(J+1)N=O(n)$. Thus maximum degree is
$\sqrt b\,\polylog b$ and parallel multiplicity is $\polylog b$,
including the factor $R$ from merging old owners. After word recoding
the root has size $b^{1+o(1)}$ with the same required degree and
multiplicity bounds up to polylogarithms. Local entry descriptions and
arithmetic cost are polynomial in $b,m$. Public blocks can be
transformed by direct finite sums in polynomial time. All bits of each
transmitted transform word share endpoints.

\section{Size and end-to-end balance}
The original $m$ long hash-output arrays contributed $(4+o(1))mb$
Boolean word edges. Remove these arrays and their old hash carries.
Keep the existing coordinate/slope inputs, quotient inputs and
arithmetic outputs. They contribute $(34+o(1))mb$ because the original
total was $(38+o(1))mb$. The replacement word cost per slot is
\[
  (J+1)NW=(2+o(1))b.
\]
Indeed, $n/H\le J<n/H+1$, $N/n=o(1)$ and $N/H=1+o(1)$.
Also $nk=b+O(k)$ and $W/k\to2$. Counting the shared short-input
transform is essential here. Thus the new total is $(36+o(1))mb$.
Rounding word widths into the fixed simplex alphabet costs $O(mb/k)$
symbol edges. Both it and the auxiliary binary costs remain lower order
after multiplication by the dimension $t=O(\log\log b)$.

With $K=2\pi\sqrt{14\log_2(90)}/(9\ln2)$ from checkpoint 004,
the inherited conditional contraction theorem gives
\[
 \log_2 C_m \le (18K+o(1))\,mb/\sqrt{\log_2 b}.
\]
The slot distribution and all accepted ordered signed tuples are
unchanged. Factor enumeration, the guarded counting and pinning
interface, sampling, training, and the verified solver therefore admit
the original total bound $\operatorname{poly}(b,m)(B_0 C_m+mB_0^2)$
with $\log_2 B_0=b/m+O(\log b+\log m)$. The sampling proof uses an
identical deterministic selection rule for training and online calls
to the new finder. Its coverage premise uses relation existence.
Choosing $m=(18K)^{-1/2}(\log_2 b)^{1/4}(1+o(1))$ gives leading
time coefficient $2\sqrt{18K}=26.2933775155911$ and space coefficient
$\sqrt{18K}=13.1466887577956$. The ratio to the old time coefficient
is $\sqrt{18/19}$. These are asymptotic conditional upper bounds.

\section{Count all tuples at one specified target}
The above graph directly counts C38's generic tuples. The ordinary
convolution coefficient $h(T)$ also includes intermediate identities
and equal-abscissa additions. The following construction obtains that
integer by one guarded graph count, including for $T=O$. It neither
builds partial-sum sets nor queries interval coefficients.

If $T\ne O$, count generic extended paths with factors
\[
 D,A_1,\ldots,A_m,-D
\]
and last partial point $T$. The original slot restrictions apply only
to the $A_i$. If $T=O$, insert a fixed nonidentity point $Q$ immediately
before $-D$ and use final target $Q$. The supplied subgroup generator
suffices for $Q$. Write $Z$ for the resulting integer graph count.

For any tuple reaching the specified target, start with all $n_E$
possible offsets $D$. Exclude $D=O$. At an interior step with fixed
addend $A$ (including the optional $Q$), the previous partial is
$D+S$ for a fixed $S$. The two equal-abscissa exceptions $D+S=\pm A$
exclude at most two offsets. The inverse case already excludes an
intermediate identity. At the closing step the previous partial is
$D+W$ with fixed nonidentity $W$ equal to the final target.
Equality with $D$ would imply $W=O$, and equality with $-D$ requires
$2D=-W$, with at most four solutions over an odd prime field.
Thus at most $d=2m+7$ offsets per tuple are lost. Every counted path
has the requested ordinary sum after cancellation. Consequently
\[
    0\le n_Eh(T)-Z\le d\,h(T).
\]
For any slot $i$, fixing the other points determines its required group
element uniquely. Hence $h(T)\le U:=\min_i\prod_{j\ne i}|F_j|
\le (2\tau)^{m-1}$. If $dU<n_E$, the exact answer is
\[
                 h(T)=\left\lceil Z/n_E\right\rceil .
\]
The coarser test $d(2\tau)^{m-1}<n_E$ uses public parameters only.
It holds uniformly eventually at the project's threshold for
polylogarithmic $m$, since its left side has logarithm at most
$b-b/m+O(m\log b)$ and $\log_2 n_E=b+O(1)$. Outside the rounding
condition, direct finite streaming of slot tuples is an exact fallback.
The rounding argument is deterministic, including when $Z=0$.

\paragraph{Pricing the offset inside the graph.}
The first and last abscissas are constrained equal at the source
owners. If the extended arity is $n'$, the two ordinates are
\[
 y_1=\lambda_2(x_1-X_2)-Y_2,\qquad
 y_{n'}=\lambda_{n'}(x_{n'}-X_{n'})-Y_{n'} .
\]
Impose $y_1+y_{n'}=0\bmod p$ as one additional quadratic equation.
For a fixed penultimate $Q$, pin its abscissa. At its index $j$, pin
its ordinate with $\lambda_j(x_j-X_j)-Y_j$ if $j$ is even and
$\lambda_j(x_j-X_{j-1})+Y_{j-1}$ if $j$ is odd. These expressions
use only retained or public ordinates. The added integer residuals
fit the original shifted quotient and carry bounds. There are a
constant number of additional slots, arrays and equations, costing
$O(b)$ word edges. The leading relation-count coefficient is unchanged
for $m\to\infty$. No external loop over the $n_E$ possible offsets is
part of the conditional algorithm. For the original discrete-log
application use the guarded count directly to select witnesses.

\paragraph{Finite evidence and scope.}
Run \texttt{python -B research/short\_hash\_count/verify.py}.
The checks exercise true forward and inverse Fourier row interfaces,
including blocked overlaps, top partial limbs, discarded overflow,
canonical input restrictions and carry mutations. There are 8,192
exhaustive four-bit cases, 48 larger block cases, and 20 rejected
corruptions in the latter suite. Two retained C38 predicate fixtures
use alphabets of sizes two and six. Their old hash-only symbols are
removed before the arithmetic and replacement hash tests run.
Tiny selected-target
fixtures compare literal ordinary counts, old guarded counts, and
rounded extended counts. The verifier enumerates offsets on the toy
instances to check multiplicities independently. For example on the
order-223 curve over $\mathbb F_{211}$, the extended count for
$(170,205)$ is 217, recovering one ordinary tuple even though the
original guarded count is zero. There is also an all-target check on
an order-64 curve with full rational two-torsion. Large ledger cases
evaluate integer formulas only. The supplied asymptotic contractor is
a premise. These runs are not end-to-end time benchmarks.
\end{document}
